\scanpage{26}
\begin{center}-2.22-\end{center}
of some torus $T\subset H$. If $\epsilon_1\neq\epsilon_2$, then there do not exist positive integers $n$ and $m$ such that $\lambda_1(\alpha^n)=\lambda_2(\alpha^m)$, all $\alpha\in\mathbb G_m$. But then there is a character $\chi$ of $T$ such that $\chi\circ\lambda_1=0$, (the trivial homomorphism from $\mathbb G_m$ to $\mathbb G_m$), and $\chi\circ\lambda_2\neq0$. Since the kernel of $\phi$ is finite, there is an integer $N$ such that $\chi^N$ is induced by a character $\chi'$ of $\phi(T)$. Now suppose that $\phi_*\epsilon_1=\phi_*\epsilon_2$. Then $P(\phi\circ\lambda_1)=P(\phi_*\epsilon_1)=P(\phi_*\epsilon_2)=P(\phi\circ\lambda_2)$. Let $T''$ be the intersec-
tion of $\phi(T)$ with the radical of $P(\phi\circ\lambda_1)$, and let $\chi''$ be the restriction of $\chi'$ to $T$. According to the structure theory of parabolic subgroups (cf.\ \hspace{2cm}), some power of $\chi''$ extends to a character $\chi'''$ of $P(\phi\circ\lambda_1)$. But since $\phi\circ\lambda_1$ and $\phi\circ\lambda_2$ are 1-parameter subgroups of $P(\phi\circ\lambda_1)$, hence of $T''$, we have $\chi'''\circ(\phi\circ\lambda_1)=0$, and $\chi'''\circ(\phi\circ\lambda_2)\neq0$. But then, referring to the definition of $\Delta(G)$, one sees that this precludes the possibility $\Delta(\phi\circ\lambda_1)=\Delta(\phi\circ\lambda_2)$. Therefore $\phi_*\epsilon_1\neq\phi_*\epsilon_2$. \uline{QED}

\uline{Corollary 2.11}: If $\phi$ is an isogeny, $\phi_*$ is an isomorphism.

\uline{Proof}: It remains to show that $\phi_*$ is surjective. But if $\lambda$ is any 1-parameter subgroup of $G$, then for some positive integer $n$, and for some 1-parameter subgroup $\mu$ of $H$, $\phi\circ\mu(\alpha)=\lambda(\alpha^n)$, $\alpha\in\mathbb G_m$. \uline{QED}

Suppose we apply Proposition 2.11 to the inclusion of a maximal torus $T$ in $G$. This induces an inclusion of $\Delta(T)$ in $\Delta(G)$. But $\Delta(T)$ is very simple - it is essentially the set of rational points on a sphere. In fact, it follows immediately from the definition that $\Delta(T)$ is the set of rays in $\unsure{\Gamma|\Pi}^{\mathbb R}(T)$ which contain a point of $\unsure{\Gamma|\Pi}^{\mathbb Q}(T)$.
\mspage{25}
